% ============================================================
%  HCMUT-DEE Thesis Kit v2.0.3 — Chương 5
%  Copyright (c) 2026 Nguyen Trong Thang & TS. Nguyen Phuc Khai
%  License: MIT
% ============================================================

\chapter{Phụ Lục, Mã Nguồn \& Listing Code}
\label{ch:code-listing}

\section{Viết phụ lục}

Phụ lục được khai báo sau lệnh \verb|\appendix| trong \texttt{thesis.tex}. Mỗi phụ lục là một \verb|\chapter{}| và được đánh ký hiệu A, B, C\ldots thay vì số.

\begin{lstlisting}[style=latexstyle, caption={Khai báo phụ lục}]
% Trong thesis.tex, sau phan Tai lieu tham khao:
\appendix
% ============================================================
%  HCMUT-DEE Thesis Kit v2.0.3 — Phụ lục mẫu
%  Copyright (c) 2026 Nguyen Trong Thang & TS. Nguyen Phuc Khai
%  License: MIT
% ============================================================

\chapter{Phụ lục mẫu}
\label{app:sample}

% ═══════════════════════════════════════════════════════════════
%  GHI CHU: Day la phu luc mau. Hay thay the bang noi dung
%  phu luc cua do an ban. Vi du: ma nguon, du lieu, bang phu.
% ═══════════════════════════════════════════════════════════════

\section{Ví dụ mã nguồn Python}

\begin{lstlisting}[style=pythonstyle, caption={Hàm tính tổn thất công suất trên lưới điện}, label={lst:app-python}]
import numpy as np

def backward_forward_sweep(R, X, P_load, Q_load, V_base=12.66):
    """
    Giai luu luong cong suat bang phuong phap
    Backward/Forward Sweep cho luoi hinh tia.
    """
    n_bus = len(P_load)
    V = np.ones(n_bus) * V_base  # Khoi tao dien ap

    for iteration in range(100):
        # Backward sweep: tinh dong nhanh
        I_branch = np.zeros(n_bus - 1, dtype=complex)
        for k in range(n_bus - 2, -1, -1):
            S = complex(P_load[k+1], Q_load[k+1])
            I_branch[k] = np.conj(S / V[k+1])

        # Forward sweep: cap nhat dien ap
        for k in range(n_bus - 1):
            Z = complex(R[k], X[k])
            V[k+1] = V[k] - I_branch[k] * Z

        # Kiem tra hoi tu
        if np.max(np.abs(V - V_old)) < 1e-6:
            break
        V_old = V.copy()

    return V
\end{lstlisting}

\section{Ví dụ bảng dữ liệu}

\begin{table}[H]
\centering
\caption{Thông số hệ thống IEEE 33-bus (trích một phần)}
\label{tab:app-ieee33}
\begin{tabular}{cccccc}
\toprule
\textbf{Nhánh} & \textbf{Từ} & \textbf{Đến} & \textbf{R ($\Omega$)} & \textbf{X ($\Omega$)} & \textbf{P$_{load}$ (kW)} \\
\midrule
1  & 1 & 2  & 0.0922 & 0.0477 & 100 \\
2  & 2 & 3  & 0.4930 & 0.2511 & 90  \\
3  & 3 & 4  & 0.3661 & 0.1864 & 120 \\
4  & 4 & 5  & 0.3811 & 0.1941 & 60  \\
5  & 5 & 6  & 0.8190 & 0.7070 & 60  \\
\bottomrule
\end{tabular}
\end{table}

% === CHEN FILE PDF BEN NGOAI (bo comment neu can) ===
% \section{Datasheet thiet bi}
% \includepdf[pages=-]{path/to/datasheet.pdf}
   % Phu luc A
% \input{appendices/appendix-b}       % Phu luc B (neu can)
\end{lstlisting}

\section{Chèn mã nguồn với listings}

Template đã định nghĩa sẵn \textbf{4 style} cho các ngôn ngữ phổ biến. Bạn chỉ cần chọn style phù hợp.

% ═══════════════════════════════════════════════════════════════
\subsection{Python (\texttt{pythonstyle})}

\begin{lstlisting}[style=pythonstyle, caption={Ví dụ mã Python --- Tính tổn thất công suất}, label={lst:python-sample}]
import numpy as np

def calculate_power_loss(V, I, R):
    """Tinh ton that cong suat tren luoi dien."""
    P_loss = np.sum(I**2 * R)
    print(f"Tong ton that: {P_loss:.2f} kW")
    return P_loss

# Tinh toan
voltage = np.array([1.0, 0.98, 0.95])
current = np.array([50.0, 45.0, 60.0])
resistance = np.array([0.01, 0.02, 0.015])
loss = calculate_power_loss(voltage, current, resistance)
\end{lstlisting}

% ═══════════════════════════════════════════════════════════════
\subsection{MATLAB (\texttt{matlabstyle})}

\begin{lstlisting}[style=matlabstyle, caption={Ví dụ mã MATLAB --- Giải hệ phương trình}, label={lst:matlab-sample}]
% Giai he phuong trinh tuyen tinh Ax = b
A = [4 -1 0; -1 4 -1; 0 -1 4];
b = [15; 10; 10];

% Phuong phap truc tiep
x = A \ b;

% Hien thi ket qua
fprintf('Nghiem x = [%.4f, %.4f, %.4f]\n', x);

% Ve do thi
figure;
plot(x, 'ro-', 'LineWidth', 2);
xlabel('Node'); ylabel('Voltage (p.u.)');
title('Voltage Profile');
grid on;
\end{lstlisting}

% ═══════════════════════════════════════════════════════════════
\subsection{C/C++ (\texttt{cstyle})}

\begin{lstlisting}[style=cstyle, caption={Ví dụ mã C --- Đọc ADC vi điều khiển}, label={lst:c-sample}]
#include <stdio.h>
#include <stdint.h>

#define ADC_RESOLUTION  4096    // 12-bit ADC
#define V_REF           3.3     // Dien ap tham chieu (V)

// Ham doc gia tri dien ap tu ADC
float read_voltage(uint16_t adc_value) {
    float voltage = (float)adc_value / ADC_RESOLUTION * V_REF;
    return voltage;
}

int main(void) {
    uint16_t raw = 2048;  // Gia tri ADC thu
    float v = read_voltage(raw);
    printf("Dien ap: %.3f V\n", v);
    return 0;
}
\end{lstlisting}

% ═══════════════════════════════════════════════════════════════
\subsection{PLC Structured Text (\texttt{plcstyle})}

\begin{lstlisting}[style=plcstyle, caption={Ví dụ PLC Structured Text --- Điều khiển động cơ}, label={lst:plc-sample}]
PROGRAM Motor_Control
VAR
    Start_Button  : BOOL := FALSE;
    Stop_Button   : BOOL := FALSE;
    Motor_Running : BOOL := FALSE;
    Speed_SP      : REAL := 1500.0;    (* Toc do dat (RPM) *)
    Speed_PV      : REAL := 0.0;       (* Toc do thuc te *)
    Timer_Delay   : TON;               (* Bo dinh thoi *)
END_VAR

(* Logic dieu khien chinh *)
IF Start_Button AND NOT Stop_Button THEN
    Timer_Delay(IN := TRUE, PT := T#2S);
    IF Timer_Delay.Q THEN
        Motor_Running := TRUE;
    END_IF;
ELSIF Stop_Button THEN
    Motor_Running := FALSE;
    Timer_Delay(IN := FALSE);
END_IF;

END_PROGRAM
\end{lstlisting}

% ═══════════════════════════════════════════════════════════════
\section{Chèn file code bên ngoài}

Thay vì copy code vào file \texttt{.tex}, bạn có thể include trực tiếp:

\begin{lstlisting}[style=latexstyle, caption={Chèn file code bên ngoài}]
% Chen file Python
\lstinputlisting[style=pythonstyle,
  caption={Ma nguon chuong trinh chinh},
  label={lst:main-py}
]{path/to/main.py}

% Chen 1 phan cua file (dong 10 den 25)
\lstinputlisting[style=matlabstyle,
  firstline=10, lastline=25,
  caption={Ham tinh luu luong cong suat}
]{path/to/powerflow.m}
\end{lstlisting}

% ═══════════════════════════════════════════════════════════════
\section{Đính kèm PDF bên ngoài}

Dùng package \texttt{pdfpages} để chèn file PDF (datasheet, chứng chỉ, \ldots):

\begin{lstlisting}[style=latexstyle, caption={Chèn file PDF bên ngoài}]
% Chen tat ca cac trang cua file PDF
\includepdf[pages=-]{path/to/document.pdf}

% Chi chen trang 1 va 3
\includepdf[pages={1,3}]{path/to/document.pdf}

% Chen voi tieu de
\includepdf[pages=-, pagecommand={
  \thispagestyle{fancy}}]{path/to/datasheet.pdf}
\end{lstlisting}
